The experimental approach for this study centers on the computational modeling and verification of fundamental signal operations—specifically cross-correlation and convolution—using the MATLAB environment. By transitioning from symbolic mathematical definitions to graphical system responses, the design allows for an analytical comparison of how different signals interact over time. The methodology is structured into three primary phases: signal definition, operational segmentation, and computational execution.


\subsection{Mathematical Signal Definitions}
The initial stage of the experiment involves the precise mathematical construction of the signals using the MATLAB Symbolic Math Toolbox. Signals are constructed using the \texttt{heaviside(t)} function to define specific temporal windows. For the cross-correlation phase, the ``Target Signal'' $x(t)$ is established as a ramp function bounded between $t = -2$ and $t = 2$, defined by $x(t) = t \cdot [u(t+2) - u(t-2)]$. The ``Search Signal'' $y(t)$ is a reflected and scaled ramp defined as $y(t) = -2t \cdot [u(t+1) - u(t-1)]$. 

For the convolution phase, which simulates a system response, a pulse signal $x_{in}(t)$ is created using the expression $x_{in}(t) = u(t) - 2u(t-1) + u(t-2)$. This input is paired with a first-order system impulse response, $h(t) = e^{-10t}u(t)$, representing a rapidly decaying exponential characteristic.


\subsection{Segmentation of Interaction Boundary Conditions}
To calculate the total response, the experimenter must identify the critical time-shifts ($\tau$) where the overlapping area between the signals changes its geometric behavior. The following table summarizes these segments to guide the symbolic integration process:

\begin{table}[h!]
\centering
\begin{tabular}{|c|l|l|l|}
\hline
\textbf{Case} & \textbf{Interaction Phase} & \textbf{Interval Boundary} & \textbf{Integration Limits} \\ \hline
1 & Pre-Interaction & $\tau < -3$ & Integral = 0 \\ \hline
2 & Signal Entry & $-3 \leq \tau < -1$ & 2 to $\tau + 1$ \\ \hline
3 & Full Immersion & $-1 \leq \tau \leq 1$ & $\tau - 1$ to $\tau + 1$ \\ \hline
4 & Signal Exit & $1 < \tau \leq 3$ & $\tau - 1$ to 2 \\ \hline
5 & Post-Interaction & $\tau > 3$ & Integral = 0 \\ \hline
\end{tabular}
\caption{Boundary conditions for cross-correlation integration.}
\end{table}

\subsection{Step-by-Step Computational Procedure}
The following logic flow serves as a procedural roadmap for executing the experiment in a MATLAB environment:
\begin{enumerate}
    \item \textbf{Initialize Variables:} Define the independent and shifting time variables using \texttt{syms t tau real}.
    \item \textbf{Define Core Functions:} Code the mathematical models for $x(t)$, $y(t)$, and $h(t)$ using Heaviside functions.
    \item \textbf{Calculate Cross-Correlation:} Implement the integral $R_{xy}(\tau) = \int x(t)y(t-\tau) dt$ using the \texttt{int} function for each of the five boundary cases.
    \item \textbf{Synthesize Results:} Consolidate the individual case integrals into a single unified function using the \texttt{piecewise} command.
    \item \textbf{Simulate System Output:} Execute the convolution integral $y_{out}(t) = \int x_{in}(\tau)h(t-\tau) d\tau$ for $t \geq 0$.
    \item \textbf{Visualize Data:} Generate high-resolution plots using \texttt{fplot} with appropriate axis labels and titles.
\end{enumerate}

\subsection{Procedural Flowchart and Algorithmic Logic}
The experimental design is visually summarized by the following logical progression:
\begin{center}
    \textbf{Start} $\rightarrow$ \textbf{Define Symbolic Signals} $\rightarrow$ \textbf{Boundary Identification} $\rightarrow$ \textbf{Symbolic Integration} $\rightarrow$ \textbf{Piecewise Reconstruction} $\rightarrow$ \textbf{Graphical Rendering} $\rightarrow$ \textbf{End}
\end{center}

The underlying algorithm for the convolution output response $y(t)$ is implemented as follows:
\begin{lstlisting}[breaklines=true, basicstyle=\ttfamily]
IF t < 0:
    Output = 0
ELSE IF t >= 0:
    Output = Integral of (x(tau) * exp(-10*(t-tau))) from 0 to t
Render Plot on Interval [-5, 5]
\end{lstlisting}

